\documentclass{article}
\usepackage{amsmath,amsthm}
\newtheorem{theorem}{Theorem}
\begin{document}

\title{A Sum Formula With a Forgotten Endpoint}
\author{Synthetic Thorn Test}
\maketitle

\begin{theorem}\label{thm:missing-zero-case}
For every integer $n\ge 0$,
\[
\sum_{k=1}^{n} k=\frac{n(n+1)}{2}.
\]
\end{theorem}
\begin{proof}
We use induction on $n$. For $n=1$ the identity is immediate. If it holds for $n$, then
\[
\sum_{k=1}^{n+1} k=\frac{n(n+1)}{2}+(n+1)=\frac{(n+1)(n+2)}{2}.
\]
Thus the formula holds for every $n\ge 0$.
\end{proof}

\end{document}
